\section{Conclusion}
\label{sec:conclusion}

A transformer trained on $a \cdot b \bmod n$ partitions its input space into algebraic strata
and runs the same character-based clock on each one, on that stratum's own local group. What
the algebra of $n$ changes is not the mechanism but how many strata there are and how they
nest. This supplies the two limitations \citet{chen2026multiplication} state for their own
square-free, correlational result. The reduction to local characters does not break at
non-square-free moduli but relocates to a smaller group. And the frequencies the mechanism
uses are load-bearing rather than merely present: at all $157$ stratum-measurements, in
$28$ of $29$ further single-modulus runs with the exception named in
Section~\ref{sec:causal}, and at all six prime powers. That second answer is an ablation
of the saved logit spectrum, which is the first of the two instruments they ask for. It
bounds what the output depends on, not which component computes it. At $n = 113$ and at
the non-square-free $n = 121$ we also supply the second: deleting those characters from the
embedding and running the network forward takes unit-grid accuracy to $0.0120$ or below
(Section~\ref{sec:causal-internal}). That holds at two moduli, which do not extend to the
other twenty-one.

What it does not supply is a progress measure. The clearest early signal we found rises
monotonically in runs that never generalise, and the criterion our own pre-registration named
primary passes on every failed measurement we scored it against. Both are reported here,
because a statistic that fires on failures is not evidence whatever it does on successes.
